\chapter{Quantitative Investment Strategies}\label{ch:s2:quantitative-investment-strategies}

A bank can package a carry or momentum strategy as an index, license it, and sell swaps on it. The backtest that sells the index is written before the
index goes live, and its returns after launch are lower. Suhonen, Lennkh and Perez studied 215 such strategies developed and promoted by global
investment banks and found a median 73\% fall in Sharpe ratios from the backtest to the live period. On this chapter's synthetic banks, each launching
the best three of twenty candidate indices, the launched indices backtest at a Sharpe ratio of 0.51 and live at 0.16 after fees, and more than a third
of them lose money live. The build is \texttt{firm.qis}.

\section{Rules-based indices as products}

\begin{definition}[Quantitative investment strategy]\label{def:s2:quantitative-investment-strategies:qis}
A \emph{quantitative investment strategy}\index{quantitative investment strategy} (QIS) is a bank's rules-based systematic strategy published as an
index (a systematic strategy index, Book~5, chapter~19), with a rulebook, a calculation agent and a backtest, which clients access through swaps,
notes or funds rather than by running it themselves.
\end{definition}

A QIS index is a strategy with a contract around it: the rules fix what the index holds, when it rebalances and what costs it deducts, so the client
can buy exposure to carry, momentum, value or a volatility premium without a team to run it. The bank earns fees and the margin on the swaps, and it
decides which indices to launch. That decision is made on backtests.

\texttt{firm.qis} gives 2,000 product teams twenty candidate indices each (\cref{lst:s2:quantitative-investment-strategies:launch}). The candidates'
true Sharpe ratios are drawn around 0.10 with a standard deviation of 0.15, and their estimation errors are correlated at 0.5, since they are variants
of a few ideas. Each team backtests over ten years, launches its best three, and lives with them for five, at a 10\% volatility target. The live
indices pay 0.5\% a year of index and swap fees and lose 0.2\% a year to predictable rebalancing. The true Sharpe ratios are calibrated so that the
median decay matches the 73\% Suhonen and co-authors report; everything else follows from selection.

\section{Wrappers and fees}

\begin{definition}[Strategy wrapper]\label{def:s2:quantitative-investment-strategies:wrapper}
A \emph{strategy wrapper}\index{strategy wrapper} is the legal and economic form through which a client holds a strategy index: a total-return swap,
a structured note, a certificate or a fund, each with its own fees, counterparty, capital treatment and liquidity.
\end{definition}

At a 10\% volatility target, 0.7\% a year of fees and leakage costs 0.07 of Sharpe ratio. That is a large share of a true Sharpe ratio of 0.1 to 0.3,
and it is invisible in a backtest run before the index existed: the backtest has no fee because nobody was paying one. The wrapper adds its own terms:
a swap carries counterparty risk and a funding spread, and a note adds the issuer's credit spread (Book~5, chapter~19).

\section{Hedging the index}

A bank that pays a client an index's return on a swap hedges by holding the index's positions. The rulebook tells everyone when the index will
rebalance and roughly what it will trade, so others can trade ahead of it, and the index buys a little higher and sells a little lower than it would
unobserved. The chapter charges this as a leakage of 0.2\% a year, borne only live, since a backtest trades at the historical close. The same
predictability lets the bank net the index's trades with its other flows (chapter~25), which lowers its own cost of hedging but not the client's.

\section{Live against backtest}

\begin{definition}[Index backtest decay]\label{def:s2:quantitative-investment-strategies:decay}
\emph{Index backtest decay}\index{index backtest decay} is the fall from the Sharpe ratio of a strategy index's pre-launch backtest to that of its
live, after-fee returns, measured across launched indices.
\end{definition}

\begin{center}
\small
\begin{tabular}{@{}lcc@{}}
\toprule
Sharpe ratios, mean over launched indices & best three of twenty & best of twenty\\
\midrule
backtest (ten years) & 0.51 & 0.62\\
true & 0.22 & 0.26\\
live, before costs (five years) & 0.23 & 0.27\\
live, after fees and leakage & 0.16 & 0.20\\
median decay, backtest to live after costs & 72.2\% & 69.4\%\\
\bottomrule
\end{tabular}
\end{center}

The live Sharpe ratio is close to the true one: live returns are honest, only noisy. The backtest is not, because it was chosen for being high. Selection
did find better strategies, with a true Sharpe ratio of 0.22 against 0.10 for the average candidate, but it credited them with noise as well. After
costs, 36.6\% of the launched indices had a negative live Sharpe ratio. One team's best index shows how it looks from inside
(\cref{fig:s2:quantitative-investment-strategies:path}): a backtest Sharpe ratio of 0.71 and a gain of 71\% over ten years, then five live years with a
Sharpe ratio of $-0.07$ before costs and a loss of 7.1\% after them.

\begin{omfigure}
\begin{tikzpicture}
\begin{axis}[width=11cm,height=5.2cm,xlabel={\small years from launch},ylabel={\small cumulative return (\%)},tick label style={font=\footnotesize},
  xmin=-10,xmax=5,ymin=-20,ymax=90,ytick={-20,0,20,40,60,80},grid=major,grid style={gray!25},table/col sep=comma]
\fill[gray!15] (axis cs:0,-20) rectangle (axis cs:5,90);
\addplot[omThm,thick] table[x=year,y=cum]{figdata/strategies-2/27-quantitative-investment-strategies/path.csv};
\node[font=\scriptsize,anchor=north west] at (axis cs:-9.6,85) {backtest};
\node[font=\scriptsize,anchor=north west] at (axis cs:0.3,85) {live};
\end{axis}
\end{tikzpicture}

\omcaption{One synthetic product team's best index: cumulative sum of daily returns at a 10\% volatility target over its ten-year backtest and its first
five live years after fees and leakage (shaded). Data: \texttt{s2\_qis.path}.}\label{fig:s2:quantitative-investment-strategies:path}
\end{omfigure}

The deflated Sharpe ratio (Book~4, chapter~12) asks whether the best of $N$ backtests beats what the best of $N$ worthless ones would show. Bailey and
Lopez de Prado built it to correct for exactly this selection. Applied to each team's best index, with twenty trials and the team's own dispersion of
backtest Sharpe ratios, it passes 5.9\% of them at 95\% confidence. Those that pass backtested at 1.13 and live at 0.32 before costs; those that fail,
at 0.59 and 0.27. The test is right to be sceptical, but ten years of backtest cannot tell a Sharpe ratio of 0.2 from one of 0.3: the discipline it buys
is in the backtests it stops from being sold as excellent, not in picking winners.

Suhonen and co-authors also found that the most complex strategies lost over 30 percentage points more than the simplest. Complexity means more
candidates tried, openly or not. In the model, choosing the best of more candidates raises the backtest far more than the live result
(\cref{fig:s2:quantitative-investment-strategies:complexity}): from 2 to 100 candidates, the best backtest rises from 0.26 to 0.76 and the live Sharpe
ratio after costs from 0.07 to 0.24, so the gap between them grows from 0.18 to 0.52. McLean and Pontiff found a related decay in published equity
predictors, whose returns were 26\% lower out of sample and 58\% lower after publication.

\begin{omfigure}
\begin{tikzpicture}
\begin{axis}[width=11cm,height=5.2cm,xmode=log,xlabel={\small candidate indices tried (log scale)},ylabel={\small Sharpe ratio of the best},
  tick label style={font=\footnotesize},xmin=2,xmax=100,ymin=0,ymax=0.85,xtick={2,5,10,20,50,100},xticklabels={2,5,10,20,50,100},grid=major,
  grid style={gray!25},table/col sep=comma,legend style={font=\scriptsize,at={(0.03,0.97)},anchor=north west}]
\addplot[omThm,thick,mark=*,mark size=1.2pt] table[x=candidates,y=backtest]{figdata/strategies-2/27-quantitative-investment-strategies/complexity.csv};
\addplot[omDef,thick,mark=square*,mark size=1.2pt] table[x=candidates,y=net]{figdata/strategies-2/27-quantitative-investment-strategies/complexity.csv};
\legend{backtest,live after costs}
\end{axis}
\end{tikzpicture}

\omcaption{The best of $N$ candidate indices: mean backtest Sharpe ratio and mean live Sharpe ratio after fees and leakage, over 2,000 synthetic
product teams, against the number of candidates tried. Data: \texttt{s2\_qis.complexity}.}\label{fig:s2:quantitative-investment-strategies:complexity}
\end{omfigure}

\section{Strategy files}

\begin{strategyfile}{Risk-premia index swap}\label{strat:s2:quantitative-investment-strategies:rp}
\sfield{Who pays you, and why} Clients who want systematic premia without running them; the bank earns fees and swap margin.
\sfield{Instruments and venues} Total-return swaps on the bank's indices; the index's underlying futures and stocks as the hedge.
\sfield{Signal} The index rulebook.
\sfield{Sizing and execution} Hedge by running the index; net its trades with other flows.
\sfield{Costs} Rebalancing leakage; funding; capital.
\sfield{How it dies} Live returns far below the backtest; clients leave.
\sfield{Horizon, capacity, infrastructure} Years; an index calculation agent and a hedging desk.
\sfield{Backtest honestly} Count every candidate tried; deduct fees and leakage; report the deflated Sharpe ratio.
\sfield{Sources} Suhonen, Lennkh and Perez (2017); this chapter: backtest 0.51, live after costs 0.16.
\end{strategyfile}

\begin{strategyfile}{Volatility-carry index}\label{strat:s2:quantitative-investment-strategies:volcarry}
\sfield{Who pays you, and why} Buyers of insurance (chapter~1); the index sells it systematically.
\sfield{Instruments and venues} Index options or variance swaps rolled by rule.
\sfield{Signal} The rulebook: tenor, strike and size.
\sfield{Sizing and execution} Fixed notional or volatility-targeted.
\sfield{Costs} Option spreads; predictable rolls traded ahead.
\sfield{How it dies} A volatility spike; crowded rolls.
\sfield{Horizon, capacity, infrastructure} Months; option market access.
\sfield{Backtest honestly} Include the crashes and the roll costs actually paid.
\sfield{Sources} Suhonen, Lennkh and Perez (2017) found equity volatility exposures among the more robust live.
\end{strategyfile}

\begin{strategyfile}{Multi-asset risk-premia basket}\label{strat:s2:quantitative-investment-strategies:basket}
\sfield{Who pays you, and why} Diversification across carry, momentum, value and volatility premia.
\sfield{Instruments and venues} Several indices combined in one swap or fund.
\sfield{Signal} Allocation rules (equal risk).
\sfield{Sizing and execution} Volatility target at the basket level.
\sfield{Costs} Each index's fees and leakage.
\sfield{How it dies} Premia that fail together; hidden correlation.
\sfield{Horizon, capacity, infrastructure} Years.
\sfield{Backtest honestly} Each component's live record, not its backtest.
\sfield{Sources} No performance figure verified.
\end{strategyfile}

\begin{strategyfile}{Defensive overlay index}\label{strat:s2:quantitative-investment-strategies:defensive}
\sfield{Who pays you, and why} Investors paying to cut drawdowns (chapter~7).
\sfield{Instruments and venues} Trend or put-based indices overlaid on equity holdings.
\sfield{Signal} The rulebook's trend or protection rule.
\sfield{Sizing and execution} Sized to the portfolio it protects.
\sfield{Costs} The carry cost of protection; fees.
\sfield{How it dies} Protection that bleeds for years before a crash it misses.
\sfield{Horizon, capacity, infrastructure} Years.
\sfield{Backtest honestly} Judge it as insurance: cost against losses avoided, live.
\sfield{Sources} No performance figure verified.
\end{strategyfile}

\section{Tutorial: written before launch}\label{tut:s2:quantitative-investment-strategies:tut}

\begin{tutorial}
\textbf{Goal.} Simulate many product teams choosing indices from backtests, charge fees and leakage, compare live with backtest, and test the choice
with the deflated Sharpe ratio. \textbf{End state:} the table and two figures.
\begin{tutsteps}
\item \textbf{Teams and launches}.
\omcode{code/firm/qis/firm_qis.py}{48}{70}{Candidates with correlated estimation errors; the best k launched and costed.}\label{lst:s2:quantitative-investment-strategies:launch}
\item \textbf{The deflated verdict}.
\omcode{code/firm/qis/firm_qis.py}{73}{90}{Each team's best against the expected best of worthless candidates.}\label{lst:s2:quantitative-investment-strategies:dsr}
\item \textbf{Run} \texttt{launched()}, \texttt{verdicts()}, \texttt{complexity()}, \texttt{path()} and \texttt{fig\_qis.py}.
\end{tutsteps}
\textbf{What to change next.} Let teams hide failed candidates (count only the survivors) and see the deflated test fooled; add a fee that depends on
the backtest; let clients leave after two bad live years.
\end{tutorial}

\section{Build: quantitative investment strategies}\label{bld:s2:quantitative-investment-strategies:qis}

\begin{build}
\bfield{Purpose} Index selection by backtest, fees and leakage, live against backtest, and the deflated Sharpe ratio's verdict.
\bfield{Interface} \texttt{QISConfig(\dots)}, \texttt{simulate\_teams(cfg)}, \texttt{launch(sim, cfg, k)}, \texttt{deflated(sim, cfg)},
\texttt{example\_path(cfg, team)}.
\bfield{Rules} Estimated Sharpe ratios are the truth plus sampling error of $1/\sqrt{T}$; fees and leakage only live; Book 4's deflated Sharpe
ratio.
\bfield{Acceptance tests} \texttt{code/firm/qis/tests/}: the sampling error's size; selection inflates the backtest, not the truth; probabilities
in range and the example path's Sharpe ratio.
\bfield{Stretch} Non-normal returns; candidates of different complexity; hidden trials.
\end{build}

\begin{omsources}
\item A.~Suhonen, M.~Lennkh and F.~Perez, ``Quantifying backtest overfitting in alternative beta strategies'', \emph{Journal of Portfolio Management}
43(2), 2017.
\item D.~H.~Bailey and M.~Lopez de Prado, ``The deflated Sharpe ratio: correcting for selection bias, backtest overfitting, and non-normality'',
\emph{Journal of Portfolio Management} 40(5), 2014.
\item R.~D.~McLean and J.~Pontiff, ``Does academic research destroy stock return predictability?'', \emph{Journal of Finance} 71(1), 2016.
\end{omsources}

\section{Exercises}

\begin{exercise}[$\star$]\label{exo:s2:quantitative-investment-strategies:1}
An index targets 10\% volatility and costs 0.7\% a year in fees and leakage. How much Sharpe ratio do the costs take?
\end{exercise}

\begin{exercise}[$\star$]\label{exo:s2:quantitative-investment-strategies:2}
A backtest over ten years shows a Sharpe ratio of 0.5. What is its standard error, for normal returns?
\end{exercise}

\begin{exercise}[$\star$]\label{exo:s2:quantitative-investment-strategies:3}
Backtest 0.51, live after costs 0.16. What is the decay?
\end{exercise}

\begin{exercise}[$\star\star$]\label{exo:s2:quantitative-investment-strategies:4}
Why is the live Sharpe ratio close to the true one while the backtest is far above it?
\end{exercise}

\begin{exercise}[$\star\star$]\label{exo:s2:quantitative-investment-strategies:5}
Why would more complex strategies decay more?
\end{exercise}

\begin{exercise}[$\star\star$]\label{exo:s2:quantitative-investment-strategies:6}
Why does a published rebalancing schedule cost the index?
\end{exercise}

\begin{exercise}[$\star\star\star$]\label{exo:s2:quantitative-investment-strategies:7}
\emph{Coding.} Run \texttt{complexity()}. How do the best backtest and the live result change from 2 to 100 candidates, and why does the median
decay hardly move?
\end{exercise}

\begin{exercise}[$\star\star\star$]\label{exo:s2:quantitative-investment-strategies:8}
\emph{Find the flaw.} ``The index passed the deflated Sharpe test, so its live Sharpe ratio will be close to its backtest.''
\end{exercise}

\section{Problem: Written Before Launch}

\begin{problem}[{Weekend problem --- quantitative investment strategies}]\label{pb:s2:quantitative-investment-strategies:1}
The chapter's synthetic product teams and the public record.

\medskip\noindent\textbf{Part I --- The product.}
\begin{enumerate}
\item Define a \omterm{def:s2:quantitative-investment-strategies:qis}{quantitative investment strategy} and a \omterm{def:s2:quantitative-investment-strategies:wrapper}{strategy wrapper}.
\item How does the bank earn from a QIS?
\item Describe the synthetic teams and candidates.
\item What is calibrated, and to what?
\end{enumerate}

\medskip\noindent\textbf{Part II --- Costs.}
\begin{enumerate}[resume]
\item What do fees and leakage cost in Sharpe ratio?
\item Why are they absent from the backtest?
\item How does the bank hedge an index swap?
\item What does predictability cost the index?
\end{enumerate}

\medskip\noindent\textbf{Part III --- Live against backtest.}
\begin{enumerate}[resume]
\item Define \omterm{def:s2:quantitative-investment-strategies:decay}{index backtest decay}.
\item Give the table of backtest, true and live Sharpe ratios.
\item What share of launched indices lose live?
\item What does the deflated Sharpe ratio say?
\end{enumerate}

\medskip\noindent\textbf{Part IV --- The verdict.}
\begin{enumerate}[resume]
\item State the \emph{named result}: the gap between the selected indices' backtest and live Sharpe ratios.
\item What did Suhonen, Lennkh and Perez find?
\item What does complexity do in the model?
\item What did McLean and Pontiff find?
\item How should a client read a QIS backtest?
\item Which strategy file did Suhonen and co-authors find most robust?
\item How does this chapter relate to Book~7's chapter on overfitting?
\item In one sentence: what does a QIS desk sell?
\end{enumerate}
\end{problem}

\section{Interview questions}

\begin{interviewq}[$\star$ \iqroles{trader}]\label{iq:s2:quantitative-investment-strategies:1}
What is a QIS index, and how does a bank hedge a swap on one?
\end{interviewq}

\begin{interviewq}[$\star\star$ \iqroles{researcher}]\label{iq:s2:quantitative-investment-strategies:2}
A product team shows you a ten-year backtest with a Sharpe ratio of 1. What do you ask?
\end{interviewq}

\begin{interviewq}[$\star\star$ \iqroles{trader}]\label{iq:s2:quantitative-investment-strategies:3}
Other traders front-run your index's monthly rebalance. What can you change in the rulebook?
\end{interviewq}

\begin{interviewq}[$\star\star$ \iqroles{risk}]\label{iq:s2:quantitative-investment-strategies:4}
What risks does a bank keep when it sells swaps on its own indices?
\end{interviewq}

\begin{interviewq}[$\star\star$ \iqroles{developer}]\label{iq:s2:quantitative-investment-strategies:5}
What must an index calculation agent's system guarantee?
\end{interviewq}

\begin{interviewq}[$\star\star\star$ \iqroles{researcher}]\label{iq:s2:quantitative-investment-strategies:6}
Candidates' estimated Sharpe ratios are the truth plus independent noise of standard deviation $s$, and the truths are equal. Show that the expected
backtest of the best of $N$ exceeds its live expectation by about $s\sqrt{2\ln N}$ for large $N$.
\end{interviewq}
